\section*{अध्याय \ref{ch:g7:volcanoes} --- ज्वालामुखी}

\begin{solution}{exo:g7:volcanoes:1}
(क) \omterm{def:g7:volcanoes:vent}{क्रेटर}; (ख) \omterm{def:g7:volcanoes:magma}{मैग्मा}; (ग) \omterm{def:g7:volcanoes:vent}{निकास-छिद्र}; (घ) \omterm{def:g4:volcanoes-earthquakes-first-look:eruption}{लावा}; (ङ) \omterm{def:g7:volcanoes:magma-chamber}{मैग्मा कक्ष}।
\end{solution}

\begin{solution}{exo:g7:volcanoes:2}
(क) बहिःस्रावी। (ख) विस्फोटक। (ग) बहिःस्रावी: \omterm{def:g4:volcanoes-earthquakes-first-look:eruption}{लावा} का फव्वारा चमकती बूँदें
बाहर फेंकता है जो वापस गिरकर बहती हैं; \omterm{def:g7:volcanoes:magma}{मैग्मा} उड़कर राख नहीं बनता।
(घ) विस्फोटक: यह एक \omterm{def:g7:volcanoes:pyroclastic-flow}{पायरोक्लास्टिक प्रवाह} है।
\end{solution}

\begin{solution}{exo:g7:volcanoes:3}
(क) सही। (ख) ग़लत: सबसे ज़्यादा मात्रा वाली \omterm{def:g4:volcanoes-earthquakes-first-look:volcano}{ज्वालामुखी} गैस जलवाष्प है; कार्बन
डाइऑक्साइड उसके बाद आती है। (ग) ग़लत: \omterm{def:g7:volcanoes:pyroclastic-flow}{पायरोक्लास्टिक प्रवाह} गरम गैस, राख और
\omterm{def:g3:stones-pebbles-sand:rock}{चट्टान} के टुकड़ों का एक हिमस्खलन है। (घ) ग़लत: वे प्रशांत \omterm{def:g7:plate-tectonics:plate}{प्लेट} के अंदर, सबसे
पास वाली सीमा से 3200~किलोमीटर से भी ज़्यादा दूर, एक \omterm{def:g7:volcanoes:hotspot}{तप्त स्थल} के ऊपर
हैं।
\end{solution}

\begin{solution}{exo:g7:volcanoes:4}
(क) लगभग 1000 से 1200~डिग्री। (ख) लगभग 600 से 800~डिग्री। (ग) \omterm{def:g7:volcanoes:vent}{निकास-छिद्र} के
पास, जहाँ \omterm{def:g4:volcanoes-earthquakes-first-look:eruption}{लावा} नारंगी (800 से 1000~डिग्री) के बजाय पीला (1000 से 1200~डिग्री)
है: \omterm{def:g7:volcanoes:vent}{निकास-छिद्र} से दूर बहते हुए \omterm{def:g4:volcanoes-earthquakes-first-look:eruption}{लावा} ठंडा होता जाता है।
\end{solution}

\begin{solution}{exo:g7:volcanoes:5}
(क) $838 + 221 = 1059$ \omterm{def:g4:volcanoes-earthquakes-first-look:volcano}{ज्वालामुखी}। (ख) $1059 \div 1214 \approx 0.872$:
लगभग \textbf{87~\%}। (ग) $69 + 18 + 12 + 0.5 = 99.5$, लगभग 100~\%: यह छोटा
अंतर हर प्रतिशत को पूर्णांकित करने से आता है।
\end{solution}

\begin{solution}{exo:g7:volcanoes:6}
(क) $564 \div 660 \approx 0.855$: लगभग \textbf{85~\%}। (ख)
$660 - 564 - 62 - 33 = 1$: एक खड़ा शंकु। (ग) खड़े शंकु चिपचिपे \omterm{def:g4:volcanoes-earthquakes-first-look:eruption}{लावा} और राख से
बनते हैं, जो \omterm{def:g7:volcanoes:eruption-types}{विस्फोटक उद्गारों} की पहचान है: अग्नि-वलय के साथ-साथ, जहाँ एक \omterm{def:g7:plate-tectonics:plate}{प्लेट}
दूसरी के नीचे धँसती है, हम माउंट सेंट~हेलेंस और पिनातूबो जैसे बहुत-से \omterm{def:g7:volcanoes:eruption-types}{विस्फोटक
उद्गारों} की आशा करते हैं।
\end{solution}

\begin{solution}{exo:g7:volcanoes:7}
(क) $1170 - 1140 = 30$~डिग्री। (ख) $30 \div 10 = 3$~डिग्री प्रति
किलोमीटर। (ग) सुरंग में \omterm{def:g4:volcanoes-earthquakes-first-look:eruption}{लावा} हवा को नहीं छूता: उसके ऊपर की सख़्त छत गर्मी को
अंदर रखती है, जैसे बर्तन का ढक्कन। खुले में उसकी सतह तुरंत हवा को गर्मी दे देती
है।
\end{solution}

\begin{solution}{exo:g7:volcanoes:8}
(क) काउआई, कम से कम $5.5 - 0.7 = 4.8$~मिलियन वर्ष से। (ख) काउआई: \omterm{def:g4:clouds-rain-snow:precipitation}{बारिश} और
लहरें उसे कहीं ज़्यादा समय से घिस रही हैं। (ग) बिग आइलैंड, जहाँ अब भी नया \omterm{def:g4:volcanoes-earthquakes-first-look:eruption}{लावा}
निकलता है। (घ) उत्तर-पश्चिम की ओर: काउआई, जो पहले \omterm{def:g7:volcanoes:hotspot}{तप्त स्थल} से दूर ले जाया गया,
अब उसके उत्तर-पश्चिम में है।
\end{solution}

\begin{solution}{exo:g7:volcanoes:9}
(क) लगभग सारे \omterm{def:g4:volcanoes-earthquakes-first-look:volcano}{ज्वालामुखी} \omterm{def:g7:plate-tectonics:plate-boundary}{प्लेट सीमाओं} पर हैं (सूची में 87~\%): किसी भी सीमा से
इतनी दूर के \omterm{def:g4:volcanoes-earthquakes-first-look:volcano}{ज्वालामुखी} एक अपवाद हैं। (ख) हवाई के नीचे \omterm{def:g6:inside-the-earth:layers}{प्रावार} में गहरे एक \omterm{def:g7:volcanoes:hotspot}{तप्त
स्थल} मिलियनों वर्षों तक गरम रहता है। वह अपने ऊपर की \omterm{def:g3:stones-pebbles-sand:rock}{चट्टान} को आंशिक रूप से
पिघलाता है। \omterm{def:g7:volcanoes:magma}{मैग्मा}, जो आसपास की ठोस \omterm{def:g3:stones-pebbles-sand:rock}{चट्टान} से हल्का है, \omterm{def:g7:plate-tectonics:plate}{प्लेट} में से ऊपर उठता है
और फूटकर निकलता है। जैसे-जैसे \omterm{def:g7:plate-tectonics:plate}{प्लेट} चलती है, \omterm{def:g7:volcanoes:hotspot}{तप्त स्थल} के ऊपर नए \omterm{def:g4:volcanoes-earthquakes-first-look:volcano}{ज्वालामुखी} बनते
हैं और पुराने दूर ले जाए जाते हैं। (ग) नारंगी बिंदु (``किसी \omterm{def:g7:plate-tectonics:plate}{प्लेट} के
अंदर'')।
\end{solution}

\begin{solution}{exo:g7:volcanoes:10}
पहला \omterm{def:g7:volcanoes:magma}{मैग्मा} ऊपर आते हुए अपनी ज़्यादातर गैस खो चुका था, खुली छोड़ी गई या खड़ी
छोड़ी गई सोडा वाटर की बोतल की तरह: जब वह सतह तक पहुँचा, तो उसे उड़ाने के लिए
कुछ नहीं बचा था, और वह रिसकर बाहर आया। 12 और 15~जून के \omterm{def:g7:volcanoes:magma}{मैग्मा} में अब भी उसकी
गैस थी, हिलाई हुई बोतल की तरह: सतह के पास दाब गिरने पर गैस ने एक साथ बुलबुले
बना दिए, जो फूले और चिपचिपे \omterm{def:g7:volcanoes:magma}{मैग्मा} को फाड़कर राख बना दिया। पहला \omterm{def:g7:volcanoes:magma}{मैग्मा} ही वह है
जिसने अपनी गैस खो दी थी।
\end{solution}

\begin{solution}{exo:g7:volcanoes:11}
(क) $10 \div 80 = 0.125$~घंटा, और $0.125 \times 60 = 7.5$~मिनट। (ख)
$10 \div 130 \approx 0.077$~घंटा, लगभग $0.077 \times 60 \approx
4.6$~मिनट, यानी लगभग 5~मिनट। (ग) परिवार को 20~मिनट चाहिए, पर प्रवाह 5 से
8~मिनट में पहुँच जाता है: प्रवाह दिखने के बाद बहुत देर हो चुकी होती है। परिवार
तभी सुरक्षित रह सकता है जब वह उसी समय निकल जाए जब वैज्ञानिक चेतावनी के संकेत
देखते हैं, \omterm{def:g4:volcanoes-earthquakes-first-look:eruption}{उद्गार} से पहले।
\end{solution}

\begin{solution}{exo:g7:volcanoes:12}
(क) $111 \times 360 = 39\,960$: लगभग \textbf{40\,000~किलोमीटर}। (ख)
$40\,000 \div 15 \approx 2667$: लगभग 2\,670~किलोमीटर प्रति दिन। (ग)
$2667 \div 24 \approx 111$: लगभग \textbf{110~किलोमीटर प्रति घंटा}। (घ) हवा
से: \omterm{def:g4:volcanoes-earthquakes-first-look:eruption}{उद्गार} के ऊपर फेंक देने के बाद बारीक राख की अपनी कोई चाल नहीं रहती; वह हवा
में ऊँचे तैरती है और हवा उसे ले जाती है, जबकि वह धीरे-धीरे नीचे
गिरती है।
\end{solution}

\begin{solution}{pb:g7:volcanoes:1}
\textbf{1.} 16 से 31~मार्च तक: $31 - 16 = 15$~दिन; फिर अप्रैल के 30~दिन और
मई के 18~दिन: $15 + 30 + 18 = 63$~दिन।

\textbf{2.} $10\,000 \div 62 \approx 161$: औसतन, प्रति दिन 160 से ज़्यादा छोटे
\omterm{def:g4:volcanoes-earthquakes-first-look:earthquake}{भूकंप}।

\textbf{3.} फूलने से पता चला कि \omterm{def:g7:volcanoes:magma}{मैग्मा} पहाड़ के अंदर ऊपर तक उठ आया था और \omterm{def:g3:stones-pebbles-sand:rock}{चट्टान}
को बाहर धकेल रहा था। इससे उत्तरी हिस्सा खड़ा और अस्थिर भी हो गया, फिसलने को
तैयार।

\textbf{4.} नीला त्रिकोण वहाँ है जहाँ एक \omterm{def:g7:plate-tectonics:plate}{प्लेट} दूसरी के नीचे धँसती है। हम
चिपचिपे, गैस से भरे \omterm{def:g7:volcanoes:magma}{मैग्मा}, \omterm{def:g7:volcanoes:eruption-types}{विस्फोटक उद्गारों} और एक खड़े शंकु (एक स्तरित
\omterm{def:g4:volcanoes-earthquakes-first-look:volcano}{ज्वालामुखी}) की आशा करते हैं --- ठीक वही जो पहली तस्वीर दिखाती है।

\textbf{5.} $240 \div 60 = 4$~किलोमीटर प्रति मिनट।

\textbf{6.} फूला हुआ हिस्सा \omterm{def:g7:volcanoes:magma}{मैग्मा} को, दाब में, अंदर रोके हुए था, बोतल के ढक्कन
की तरह। जब वह फिसलकर हट गया, तो दाब तुरंत गिर गया; \omterm{def:g7:volcanoes:magma}{मैग्मा} में घुली गैस एक साथ
बुलबुलों में बदल गई और खुले हिस्से से फूट पड़ी --- एक तरफ़ को, जहाँ से ढक्कन हट
गया था।

\textbf{7.} $27 \div 480 \approx 0.056$~घंटा, और $0.056 \times 60 \approx
3.4$~मिनट। क्योंकि विस्फोट कम से कम इतनी तेज़ी से चला, उसे ज़्यादा से ज़्यादा
लगभग 3.4~मिनट लगे।

\textbf{8.} $8 \div 130 \approx 0.062$~घंटा, और $0.062 \times 60 \approx
3.7$: लगभग \textbf{4~मिनट}।

\textbf{9.} $24 \div 8.8 \approx 2.7$: स्तंभ माउंट एवरेस्ट से लगभग 2.7 गुना
ऊँचा उठा।

\textbf{10.} बहिःस्रावी: चिपचिपा \omterm{def:g4:volcanoes-earthquakes-first-look:eruption}{लावा} बिना फटे बाहर आया और \omterm{def:g7:volcanoes:vent}{निकास-छिद्र} के ऊपर
जमा हो गया। खड़ी छोड़ी गई बोतल की तरह, यह \omterm{def:g7:volcanoes:magma}{मैग्मा} अपनी ज़्यादातर गैस पहले ही खो
चुका था।

\textbf{11.} पहले पहाड़ हरे जंगलों के ऊपर बर्फ़ से ढका एक चिकना, नुकीला शंकु था।
बाद में उसकी चोटी और उत्तरी हिस्सा ग़ायब हैं: उनकी जगह हमारी ओर खुला एक विशाल
\omterm{def:g7:volcanoes:vent}{क्रेटर} है, और ढलानें और सामने की \omterm{def:g1:landscapes-around-us:landform}{घाटी} राख और टूटी \omterm{def:g3:stones-pebbles-sand:rock}{चट्टान} से धूसर
हैं।

\textbf{12.} $2950 - 2549 = 401$: पहाड़ ने लगभग \textbf{401~मीटर} ऊँचाई खोई,
यानी लगभग 400~मीटर, जैसा कक्षा~4 में बताया गया था।
\end{solution}
